\chapter{Plugin（插件）}

\vspace{0.5em}
\noindent\textit{前几章我们了解了 OpenClaw 的基础架构（Gateway、Agent、Channel、Binding、Session）。
本章开始进入扩展能力的篇章。Plugin（插件）是 OpenClaw 的扩展模块，
能注册新工具、新渠道、生命周期钩子等，是 OpenClaw 实现高度可扩展性的关键。}
\vspace{0.5em}

\section{概述}

Plugin（插件）是 OpenClaw 的扩展模块，可以注册：

\begin{itemize}
  \item \textbf{智能体工具}（JSON 模式函数）
  \item \textbf{渠道适配器}（新的聊天平台）
  \item \textbf{生命周期钩子}（拦截消息/工具/会话事件）
  \item \textbf{记忆提供者}
\end{itemize}

配置路径：\texttt{plugins.entries.*}

\section{安装插件}

\begin{lstlisting}[language=bash]
# From npm (recommended)
openclaw plugins install @openclaw/voice-call

# From local folder (dev)
openclaw plugins install ./extensions/my-plugin
\end{lstlisting}

安装后需\textbf{重启 Gateway 网关}加载插件。

\section{插件清单}

每个插件必须提供 \texttt{openclaw.plugin.json}：

\begin{lstlisting}[language=json]
{
  "id": "voice-call",
  "configSchema": {
    "type": "object",
    "additionalProperties": false,
    "properties": {}
  }
}
\end{lstlisting}

可选字段：\texttt{kind}、\texttt{channels}、\texttt{providers}、\texttt{skills}、\texttt{name}、\texttt{description}。

\section{注册智能体工具}

\begin{lstlisting}[language=typescript]
import { Type } from "@sinclair/typebox";

export default function (api) {
  api.registerTool({
    name: "my_tool",
    description: "Do a thing",
    parameters: Type.Object({
      input: Type.String(),
    }),
    async execute(_id, params) {
      return { content: [{ type: "text", text: params.input }] };
    },
  });
}
\end{lstlisting}

\subsection{可选工具}

可选工具需显式添加到智能体允许列表：

\begin{lstlisting}[language=json]
{
  agents: {
    list: [{
      id: "main",
      tools: {
        allow: ["workflow_tool", "group:plugins"]
      }
    }]
  }
}
\end{lstlisting}

\section{插件钩子}

插件可以注册以下生命周期钩子：

\begin{center}
\begin{tabular}{ll}
\toprule
\textbf{钩子} & \textbf{时机} \\
\midrule
\texttt{before\_agent\_start} & 运行开始前，注入上下文 \\
\texttt{agent\_end} & 完成后，检查结果 \\
\texttt{before\_tool\_call} & 工具调用前，拦截参数 \\
\texttt{after\_tool\_call} & 工具调用后，拦截结果 \\
\texttt{message\_received} & 入站消息 \\
\texttt{message\_sending} & 出站消息（发送前） \\
\texttt{session\_start / session\_end} & 会话生命周期 \\
\texttt{gateway\_start / gateway\_stop} & Gateway 生命周期 \\
\bottomrule
\end{tabular}
\end{center}

\section{配置示例}

\begin{lstlisting}[language=json]
{
  plugins: {
    entries: {
      "voice-call": {
        enabled: true,
        config: {
          provider: "twilio",
          fromNumber: "+15550001234",
          twilio: {
            accountSid: "ACxxxxxxxx",
            authToken: "..."
          }
        }
      }
    }
  }
}
\end{lstlisting}

\section{社区插件}

可通过 npm 安装社区维护的插件，例如：

\begin{itemize}
  \item \textbf{WeChat}：\texttt{openclaw plugins install @icesword760/openclaw-wechat}
  \item \textbf{Voice Call}：\texttt{openclaw plugins install @openclaw/voice-call}
  \item \textbf{Zalo}：\texttt{openclaw plugins install @openclaw/zalouser}
\end{itemize}

\section{本章小结}

本章介绍了 Plugin 插件的核心机制：

\begin{itemize}
  \item Plugin 是 OpenClaw 的\textbf{扩展模块}，可注册工具、渠道、钩子和记忆提供者
  \item 通过 \texttt{openclaw.plugin.json} 清单声明插件能力
  \item 提供丰富的\textbf{生命周期钩子}，可在消息、工具调用、会话等关键节点进行拦截
  \item 社区插件生态丰富，支持 npm 安装
\end{itemize}

\noindent 插件提供了底层的扩展能力，而更上层的“教智能体做事”的能力则由 Skill（技能）提供。下一章我们将深入了解 Skill 的工作机制。
